\section{Connes's cutoff and the main statements}
\label{sec:introduction}
Connes proved the trace formula in positive characteristic in
\citet[Section VIII, Theorem 5, p.~42]{connes1999}. The computation
at its heart is Lemma 3, pp.~44--45: without a hypothesis on zero
location, the defect of the simultaneous time and frequency cutoff
converges to the sum of the harmonic measures of all zeros, with
multiplicity. A zero on the critical line gives a point mass there.
After equation (33), pp.~46--47, Connes specifies the even prolate
selection over $\Q$ and asserts the corresponding zeta analogue.
The purpose here is to prove that growing-projection computation
for every primitive Dirichlet character and a corresponding formula
for every fixed-threshold Hankel selection at the complex place of
$\Q(\sqrt{-3})$.

Several characteristic-zero trace formulas are already theorems.
Meyer's nuclear representation gives the explicit formula with
ordinary evaluation at all zeros
\citep[Theorems~5.8 and 5.11, pp.~43, 47]{meyer2005}.
Connes and Consani give an exact fixed-scale Sonin trace identity
and an archimedean positivity estimate for restricted supports
\citep[Theorems~4.7 and 6.11, pp.~27, 48]{cc2020}.
Connes, Consani, and Marcolli give the cohomological trace and the
Weil positivity criterion
\citep[Theorem~6.1, Proposition~6.2, Corollary~6.3,
pp.~28--29]{ccm2007}.
These statements use different spaces or cutoffs. The theorem below
concerns Connes's growing finite prolate selection, its full theta
Gram matrix, every primitive conductor, and the complete Bruhat
test algebra. It gives the harmonic-measure limit before applying
any theorem locating the zeros on the line.

\subsection{The cutoff, the tests, and the zero measure}
Fix a primitive Dirichlet character $\chi$ of conductor $q$.
The raw space is the parity-$\chi(-1)$ subspace of $L^2(\R,du)$,
with its squared norm divided by two, equivalently the positive-half-line
norm. Use the raw Fourier transform with phase $e^{-2\pi iu\xi}$.
For $\Lambda>1$ let $P_\Lambda$ restrict raw time to
$[-\Lambda,\Lambda]$, and let $B_{\Lambda/q}$ restrict raw
Fourier frequency to $[-\Lambda/q,\Lambda/q]$.
The eigenfunctions of
$P_\Lambda B_{\Lambda/q}P_\Lambda$ are indexed in the full
prolate ordering, beginning at zero. Let $\Pi_{\chi,\Lambda}$
select the required parity among the indices
$n\leq\lfloor4\Lambda^2/q\rfloor$.

Put $L=\log\Lambda$, $I\psi=\int_0^\infty\psi(u)du$,
$\operatorname{ev}_0\psi=\psi(0)$, and
\[
 E_\chi\psi(t)=e^{t/2}\sum_{n\geq1}\chi(n)\psi(ne^t)
                  -\mathbf1_{\chi=1}I\psi\,e^{-t/2}.
\]
For $\chi\ne1$, let $Q_{\chi,\Lambda}$ be the orthogonal
projection in $L^2(\R,dt)$ onto
$E_\chi\operatorname{Ran}\Pi_{\chi,\Lambda}$.
For $\chi=1$ use the constrained raw space
\[
 \mathcal R_\Lambda^0=\operatorname{Ran}\Pi_{1,\Lambda}
                  \cap\ker I\cap\ker\operatorname{ev}_0
\]
to define $Q_{1,\Lambda}$. The full principal space is
$L^2(\R)\oplus\C^2$, with map
$J\psi=(E_1\psi,\psi(0),I\psi)$; its finite projection
$\mathcal Q_{1,\Lambda}$ is onto
$J\operatorname{Ran}\Pi_{1,\Lambda}$.
In nonprincipal sectors put $\mathcal Q_{\chi,\Lambda}=Q_{\chi,\Lambda}$.
All image projections use their full Hilbert-space Gram matrices.

For $a>0$ put $(V(a)f)(t)=f(t-\log a)$ and
$R(a)=V(a)\oplus a^{-1/2}\oplus a^{1/2}$ in the full principal
space. For $h\in C_c^\infty(\R)$, the space of smooth compactly
supported functions, define
\[
 V(h)=\int_\R h(r)V(e^r)dr,\qquad
 R(h)=\int_\R h(r)R(e^r)dr,\qquad
 H(z)=\int_\R h(t)e^{zt}dt.
\]
Thus $V(h)$ is logarithmic convolution and the last two components of
$R(h)$ are $H(-1/2)$ and $H(1/2)$. Our logarithmic Fourier convention
is $\widehat f(\tau)=\int f(t)e^{i\tau t}dt$, with Plancherel
measure $d\tau/(2\pi)$. Write $g^*(t)=\overline{g(-t)}$, so a
convolution square is $g*g^*$. For vectors $\|\cdot\|_2$ is the
$L^2$ norm; for operators it is the Hilbert--Schmidt norm, while
$\|\cdot\|_1$ is the trace norm. If an operator has target Hilbert
space $X$, we also write $\|T\|_{2,X}$ for its Hilbert--Schmidt
norm with that target norm. A nontrivial zero is written
$\rho=\tfrac12+\alpha_\rho+i\gamma_\rho$, with multiplicity
$m_\rho$. Define
\begin{equation}\label{eq:intro-harmonic-measure}
 d\mu_{\alpha,\gamma}(\tau)=
 \begin{cases}
 \displaystyle\frac{|\alpha|}{\pi((\tau-\gamma)^2+\alpha^2)}d\tau,
                                      &\alpha\ne0,\\
 \delta_\gamma(d\tau),&\alpha=0,
 \end{cases}
 \qquad
 \mathfrak B_\chi(h)=\sum_\rho m_\rho
                 \int_\R H(i\tau)d\mu_{\alpha_\rho,\gamma_\rho}(\tau).
\end{equation}
The sum includes every zero in $0<\Re\rho<1$ and converges
absolutely. Its point-evaluation counterpart is
$\mathfrak W_\chi(h)=\sum_\rho m_\rho H(\rho-1/2)$.

For an orthonormal basis $z_1,\ldots,z_d$ of the image defining
$Q_{\chi,\Lambda}$, let $S_L=1_{[-L,L]}$ and define its defect measure by
\begin{equation}\label{eq:intro-defect-measure}
 d\mu_{\chi,\Lambda}(\tau)=
 \left(2L-\sum_{j=1}^d|\widehat{S_Lz_j}(\tau)|^2\right)
                                                    \frac{d\tau}{2\pi}.
\end{equation}
This is independent of the basis. Bessel's inequality makes its density
nonnegative and at most $L/\pi$. When the character is fixed we write
$\mu_\Lambda$ for this measure. A Schwartz test is a smooth function
whose derivatives decrease faster than every inverse polynomial.

\begin{theorem}[The Dirichlet prolate limit, after Connes's Lemma 3]
\label{thm:main-dirichlet}
For every fixed primitive Dirichlet character and every
$h\in C_c^\infty(\R)$,
\begin{equation}\label{eq:main-sector-limit}
 \lim_{\Lambda\to\infty}
 \bigl(2\log\Lambda\,h(0)-\Tr(Q_{\chi,\Lambda}V(h))\bigr)
                  =\mathfrak B_\chi(h).
\end{equation}
There is no hypothesis of RH, simplicity, or separation of zeros.
For the full principal projection,
\begin{equation}\label{eq:main-principal-limit}
 \lim_{\Lambda\to\infty}
 \left(2\log\Lambda\,h(0)
       -\Tr(\mathcal Q_{1,\Lambda}R(h))
       +\int_\R h(t)(e^{-t/2}+e^{t/2})dt\right)
                         =\mathfrak B_1(h),
\end{equation}
with the integrated action defined above.
More strongly, the positive Fourier-diagonal measures of
$S_L-S_LQ_{\chi,\Lambda}S_L$, $S_L=1_{[-L,L]}$, converge to
$\sum_\rho m_\rho\mu_{\alpha_\rho,\gamma_\rho}$ against every
Schwartz frequency test. They have a common polynomial local-mass
bound, with constants allowed to depend on the fixed character.
\end{theorem}

The off-line weight is a probability measure, not evaluation at a
nonreal spectral parameter. The distinction is visible even for
one zero. If $\rho=1/2+\alpha+i\gamma$ with $\alpha\ne0$, the
cutoff reads the boundary average of $H$ with width $|\alpha|$.
When $\alpha=0$ that average becomes $H(i\gamma)$.
Multiplicity multiplies the measure; it does not require simple
zeros or a faithful unitary realization of a nontrivial jet block.

Identify the rational idele class group as
$C_\Q\simeq\widehat\Z^{\times}\times\R_+^{\times}$ and normalize
the first factor to have Haar measure one. Its Bruhat test algebra
consists of finite sums
\[
 h(u,e^t)=\sum_{\chi\in F}h_\chi(t)\overline{\chi(u)},
                  \qquad h_\chi\in C_c^\infty(\R).
\]
The character labels are chosen so that the $\chi$-sector action
is $\chi(u)V(e^t)$. Set
$\mathcal Q_\Lambda=\bigoplus_\chi\mathcal Q_{\chi,\Lambda}$.
Although this direct sum need not have finite rank, its product
with a Bruhat test has only finitely many finite-rank character
blocks. Define
\[
 T_\Lambda(h)=2\log\Lambda\,h(1,1)
       -\Tr(\mathcal Q_\Lambda R(h))
       +\int_\R h_1(t)(e^{-t/2}+e^{t/2})dt.
\]

\begin{corollary}[The rational trace formula and Weil's criterion]
\label{cor:main-adelic}
For every Bruhat test, with no zero-location hypothesis,
\[
             \lim_{\Lambda\to\infty}T_\Lambda(h)
                      =\sum_{\chi\in F}\mathfrak B_\chi(h_\chi).
\]
This limit equals Weil's distribution for every Bruhat test if
and only if RH holds for zeta and every primitive Dirichlet
$L$-function. Thus Alp\"oge's RH and Dirichlet GRH theorem gives
the trace formula for these projections.
\end{corollary}
The equivalence is Connes's statement. Its proof here is the
finite-character sum of Theorem~\ref{thm:main-dirichlet}, followed
by Weil's positivity criterion. The externally checked theorem
used in the last sentence concerns Dirichlet functions; its scope
does not include arbitrary Hecke or Artin functions
\citep{alpogeverification,mitverification}.

\begin{theorem}[Band comparison and the sign correction]
\label{thm:intro-band}
For $h\in C_c^\infty(\R)$ and a zero divisor satisfying
$Z_2=\sum_\rho m_\rho/(1+\gamma_\rho^2)<\infty$ and
$|\alpha_\rho|\leq\delta$, define
\[
 \mathcal N_\delta(h)=|h(0)|+
 \int_\R e^{\delta|t|}
 (|t||h''(t)|+2|h'(t)|+(1+\delta^2)|t||h(t)|)dt.
\]
Then
\[
 |\mathfrak B(h)-\mathfrak W(h)|
 \leq2\mathcal N_\delta(h)
       \sum_\rho\frac{m_\rho|\alpha_\rho|}{1+\gamma_\rho^2}
 \leq2\delta\mathcal N_\delta(h)Z_2.
\]
The distribution $\mathfrak B$ is positive on convolution squares.
Its difference from $\mathfrak W$ has no universal per-pair sign:
the formal contribution of one reflected off-line pair takes both
signs on such squares.
\end{theorem}
Theorem~\ref{thm:band-bound} and Proposition~\ref{prop:sign} prove
these assertions. The correction to the sign statement is substantive:
the converse to the trace formula uses positivity of the harmonic
distribution together with Weil's criterion, not a sign for its
difference from Weil's distribution.

\subsection{The finite-order Hecke theorem}
Let $K=\Q(\sqrt{-3})$, and let $\chi$ be a primitive finite-order
Hecke character of conductor $\mathfrak f_\chi$. Here
$N\mathfrak a=|\mathcal O_K/\mathfrak a|$ is the absolute ideal
norm, and $J_0$ is the Bessel function of the first kind of order zero. On
$L^2((0,\infty),du)$ put
\[
 c_\chi=\frac{2\pi}{\sqrt{3N\mathfrak f_\chi}},\qquad
 (\mathcal H_{c_\chi}\psi)(v)
 =c_\chi\int_0^\infty J_0(2c_\chi\sqrt{uv})\psi(u)du,
 \qquad P_\Lambda=1_{(0,\Lambda)}.
\]
Fix a concentration threshold $0<\nu<1$ and define the Hankel
prolate selection by
\[
 \Pi_{\chi,\Lambda}^{K}
 =1_{[\nu,1]}
   (P_\Lambda\mathcal H_{c_\chi}P_\Lambda
                     \mathcal H_{c_\chi}P_\Lambda).
\]
This is an explicit growing concentration subspace, independent
of the zero divisor.
Its rank is asymptotic to $2c_\chi\Lambda/\pi$.
The usual threshold $\nu=1/2$ is included.
Center the ideal theta map as
\[
 E_\chi^K\psi(t)=e^{t/2}\sum_{\mathfrak a\subset\mathcal O_K,
                                           \ \mathfrak a\ne0}
                  \chi(\mathfrak a)\psi(N\mathfrak a\,e^t)
          -\mathbf1_{\chi=1}r_KI\psi\,e^{-t/2},
 \qquad r_K=\frac{\pi}{3\sqrt3}.
\]
The character is extended by zero on ideals not coprime to its conductor.
Let $Q^K_{\chi,\Lambda}$ be its orthogonal image projection, using
the full Gram matrix and imposing $I\psi=\psi(0)=0$ in the principal
sector. Define the full principal projection by
$J_K\psi=(E_1^K\psi,\psi(0),I\psi)$ as before.
In $\mathfrak B^K_\chi$ and $\mathfrak W^K_\chi$, the sums in
\eqref{eq:intro-harmonic-measure} are now over the nontrivial zeros
of $L_K(s,\chi)$, with multiplicity.

\begin{theorem}[The finite-order Hankel limit, after Connes's Lemma 3]
\label{thm:main-hecke}
For every fixed primitive finite-order Hecke character of
$K=\Q(\sqrt{-3})$, every fixed $0<\nu<1$, and every
$h\in C_c^\infty(\R)$,
\[
 \lim_{\Lambda\to\infty}
 \bigl(2\log\Lambda\,h(0)-\Tr(Q^K_{\chi,\Lambda}V(h))\bigr)
                       =\mathfrak B^K_\chi(h).
\]
The positive Fourier-diagonal measures of
$S_L-S_LQ^K_{\chi,\Lambda}S_L$ converge to the harmonic zero measure
against every Schwartz test, with a uniform polynomial local-mass bound.
The full principal projection satisfies the same formula with the polar
correction $H(-1/2)+H(1/2)$ as in
\eqref{eq:main-principal-limit}. No hypothesis on zero location,
simplicity, or separation is imposed.

If a zero-free half-plane $\Re s>\theta$, $1/2\leq\theta<1$, is
assumed for this character and $\delta=\theta-1/2$, then
\begin{align*}
 &\left|\lim_{\Lambda\to\infty}
   \bigl(2\log\Lambda\,h(0)-\Tr(Q^K_{\chi,\Lambda}V(h))\bigr)
                  -\mathfrak W^K_\chi(h)\right|\\
 &\hspace{35mm}\leq2\delta\mathcal N_\delta(h)Z_{2,\chi},
 \qquad
 Z_{2,\chi}=\sum_\rho\frac{m_\rho}{1+(\Im\rho)^2}
                      \ll_K\log(2+3N\mathfrak f_\chi).
\end{align*}
The sector identities sum over every finite set of finite-order
characters, on the test algebra defined in
Section~\ref{sec:complex-place}.
\end{theorem}

Thus Liu's half-plane gives the explicit band width
$\delta\approx0.373332637445905$, while Hu's bound gives
$\delta=3/8$. These are applications of the cited half-plane theorems;
the cutoff limit itself uses neither. The conductor factor records the
number of zeros in the comparison. The theorem applies in particular
to the cubic Kummer characters of $K(\sqrt[3]{269})/K$, whose RH remains
equivalent to equality with Weil's distribution for all these tests.
The theorem concerns finite-order characters; it makes no assertion
for a nontrivial infinity type or for a conductor growing with the cutoff.

\subsection{Holomorphic forms and odd Artin representations}
Theorem~\ref{hnew:theorem} proves the same harmonic-measure limit
for each primitive holomorphic cuspidal newform of integer weight
$k\ge1$, level $N$, and arbitrary nebentypus. On $L^2((0,\infty),du)$,
let $\mathcal H_f$ be the transform with kernel
\[
 \frac{2\pi}{\sqrt N}
 J_{k-1}\left(\frac{4\pi\sqrt{uv}}{\sqrt N}\right),
\]
and its scalar theta map uses
$\lambda_f(n)=a_f(n)n^{-(k-1)/2}$ when
$f(z)=\sum_{n\ge1}a_f(n)e^{2\pi inz}$.
For every fixed concentration threshold the raw rank is
$4\Lambda/\sqrt N+O(\log\Lambda)$, while the logarithmic defect
still has leading term $2\log\Lambda$. The orthogonal image uses
the complete arithmetic Gram matrix. More explicitly, put
\[
 (E_f\psi)(t)=e^{t/2}\sum_{n\ge1}\lambda_f(n)\psi(ne^t),
 \qquad P_\Lambda=1_{(0,\Lambda)},
\]
\[
 \Pi_{f,\Lambda}^{\vartheta}
 =1_{[\vartheta,1]}(P_\Lambda\mathcal H_fP_\Lambda\mathcal H_fP_\Lambda),
 \qquad Q_{f,\Lambda}^{\vartheta}
 =P_{E_f\operatorname{Ran}\Pi_{f,\Lambda}^{\vartheta}}.
\]
The theta domain and positivity of the full Gram matrix are established
in Section~\ref{hnew:section}. Define $\mathfrak b_f$ and
$\mathfrak B_f(h)$ by the all-zero rule in \eqref{eq:intro-harmonic-measure},
using the nontrivial zeros of $L_f(s)=\sum\lambda_f(n)n^{-s}$.

\begin{theorem}[The holomorphic cutoff limit]\label{thm:main-holomorphic}
For every fixed primitive holomorphic cuspidal newform $f$ as above
and every fixed $0<\vartheta<1$, the positive Fourier-diagonal measures
of $S_L-S_LQ_{f,\Lambda}^{\vartheta}S_L$, $L=\log\Lambda$,
converge to $\mathfrak b_f$ against every Schwartz frequency test.
Their mass on $[T,T+1]$ is at most
$C_{f,\vartheta}(2+|T|)^{A_f}$, uniformly for $\Lambda\ge2$.
For every $h\in C_c^\infty(\mathbb R)$,
\[
 2\log\Lambda\,h(0)
   -\operatorname{Tr}(Q_{f,\Lambda}^{\vartheta}V(h))
          \longrightarrow\mathfrak B_f(h).
\]
These measure and trace conclusions also hold for the auxiliary
threshold $1-\Lambda^{-56}$. No assumption on zero location, multiplicity,
or separation is imposed, and no polar coordinates are added.
\end{theorem}

By the established weight-one modularity theorem of
\citet[Corollary~10.2(ii), author-preprint p.~21]{kharewintenberger2009}
and the all-prime conductor and Euler-factor identification of
\citet[Theorem~4.6, p.~514]{deligneserre1974}, this gives a cutoff
for every odd irreducible two-dimensional Artin representation over
$\Q$ (Corollary~\ref{hnew:odd-artin}). Both this specialization and
the holomorphic theorem are unconditional; they do not locate the
zeros on the critical line. The fixed Eisenstein-field induction is
an explicit subcase. Even degree-two and general degree-three
Artin functions are outside this conclusion.

\subsection{Proof overview}
The prolate boundary is a genuine difficulty. Landau's transition
law puts the concentration eigenvalues near $1/2$ at the full
index threshold \citep[pp.~209, 211]{landau1985}. Thus the selected
raw functions cannot be treated as having exact joint time and
frequency support. Their theta images must also be normalized by
an arithmetic Gram matrix, whose least eigenvalue is not controlled
by concentration alone.

The proof separates these issues. A fixed-character M\"obius
estimate gives an unconditional inverse-Gram bound with an
arbitrarily strong logarithmic saving. Exact theta--Poisson
identities then make each finite zero-jet packet asymptotically
orthogonal to the image. The limiting packets at the two endpoints
are Hardy model spaces; their joint Blaschke diagonal is the sum
of the harmonic measures. Same-side packets need not become
individually orthogonal, so a joint Gram calculation is essential.

For the reverse local inequality, divide a compact frequency
multiplier by $L(1/2+i\tau,\chi)$, remove its real pole jets, and approximate the
remaining raw interior by the prolate selection. The two boundary
terms converge in Hilbert--Schmidt norm. Burnol's Hardy closure
identifies precisely the parts of those boundaries that the theta
image absorbs and the zero-model parts that remain.

Passing from one fixed frequency window to all Schwartz tests
requires a further estimate. In a moving window we remove all
nearby zero jets, including off-line ones. Suitable zero-free
vertical strips give a rectangle on which the reciprocal of the $L$-function
is bounded by a fixed power of the height. A single contour
integral controls the combined residue packet, even when individual
residues become large. The zero constraints cancel its long
one-sided tails, leaving a bounded boundary triangle. The removed
packet has only logarithmic rank. This yields the polynomial
local-mass bound without a zero-spacing assumption. Finally a
smoothed omitted-tail estimate returns from the compressed
contraction to the orthogonal projection trace and recovers the
two principal polar characters.

At the complex place the raw transform is radial Fourier transformation,
expressed as a Hankel transform in the norm variable. First use the
auxiliary modes whose concentration error is at most $\Lambda^{-56}$.
They still have the full leading rank. Absolute ideal M\"obius inversion
and a Sobolev interpolation absorb the inverse-Gram loss. The endpoint
and moving-frequency arguments then give their complete trace limit.
The cross-boundary estimate uses the radial differential operator
$-\partial_u(u\partial_u)$, diagonalized by the Hankel transform.

Every fixed-threshold selection contains those auxiliary modes, so
the same harmonic measure bounds its compressed defect from above
and controls its frequency tails. For the lower inequality and the
uncompressed trace, we estimate each exterior raw functional in the
full theta norm. The exact Bessel projection kernel gives its interior
amplitude and derivative bounds. Fixed-character M\"obius cancellation
then gives an arbitrarily strong logarithmic saving for its arithmetic
dual representative. The classical mean-square theorem of
\citet[Theorem~1, printed p.~62]{chandrasekharannar1964}
controls the whole smoothed theta tail, and the same dual estimate
annihilates every fixed zero-jet packet. This proves the limit at
every fixed concentration threshold, including $1/2$.

For holomorphic forms, angular Fourier transformation replaces
the radial transform. A weighted tent gives the required Riesz
estimate with the term $L_f(0)X$ included. In weight one the
unweighted tent applies and $L_f(0)=0$. The fixed-form coefficient,
zero-free-region, and mean-square estimates then supply the same
Gram and exterior-functional bounds. Complex coefficients require
the dual theta map at the opposite endpoint; the functional
equation cancels its gamma quotient exactly. All constants keep
the form fixed, and the proof makes no separation or simplicity
assumption on its zeros.

\subsection{Comparison with the known cutoffs}
The following distinctions specify the role of the earlier work.
All page locators refer to the cited arXiv versions unless a printed
page is explicitly named.

\begin{description}[style=nextline,leftmargin=1em,font=\normalfont\bfseries]
\item[Positive characteristic and the semilocal formula.]
Connes's Lemma 3 has only finitely many zeros per character modulo
the imaginary period. The corresponding joint model-space
calculation therefore ends after a finite sum. In characteristic
zero the infinite frequency tail is an extra assertion.
Connes's Theorem 4, p.~31, concerns a fixed finite set of places
and the product $\widehat P_\Lambda P_\Lambda$, generally not an
orthogonal projection. Appendix II, Theorem 6, pp.~77--78, is the
unconditional explicit formula over every global field. It does
not by itself compute a Hilbert projection limit.

\item[The archimedean Sonin identity.]
Theorem 4.7 of \citet{cc2020} uses the even Sonin space at fixed
scale one and an explicit prolate correction distribution. Their
Theorem 6.11 restricts the support of the test to
$[2^{-1/2},2^{1/2}]$ and includes specified Mellin conditions and
an explicit correction. Appendix C, Proposition C.1, p.~51,
permits finitely many further Mellin zeros in Weil's criterion,
while still allowing arbitrary compact supports. Neither assertion
allows the global criterion to be replaced by one short interval.

\item[Hardy closure and nuclear realizations.]
Burnol's Theorem 3.2, pp.~14--15, identifies the one-sided closure
with an inner multiple of a Hardy space \citep{burnol2000}.
That theorem supplies the limiting boundary space, not a finite
prolate approximation rate. Meyer's Theorems 5.8 and 5.11 give a
summable virtual representation with the ordinary zero and pole
divisor \citep{meyer2005}. Ordinary evaluation at off-line zeros
and their positive harmonic boundary measures are different
functionals.

\item[The subsequent prolate and semilocal constructions.]
The cycle correspondence of \citet[Theorems~1.1 and 6.4,
pp.~4, 52--53]{cc2021} detects critical-line ordinates.
The self-adjoint prolate construction of
\citet[Theorems~2.6 and 6.1, pp.~7, 13]{cm2021} gives a spectral
counting asymptotic, not identification of every eigenvalue.
Theorem 4.6 of \citet[p.~23]{ccm2023} gives a bounded invertible
map between fixed-semilocal Sonin spaces. Its metric depends on
the set of places; it is not a uniform statement for a growing
set. The later ground-state constructions
\citep[Theorem~5.10, p.~23]{ccm2025} and the conditional
positivity estimate of \citet[Theorem~8.14, p.~47]{acl2026}
have hypotheses discussed below. A trace limit for fixed tests
does not supply their missing quantitative ground-vector estimate.

\item[Interpolation from zero jets.]
\citet[Theorem~1.1, p.~408; Theorems~5.1--5.2, pp.~440, 442]{bondarenko2023}
construct interpolation bases for zeta and primitive Dirichlet
functions from zero jets and logarithmic integer samples, for analytic
functions in specified strips with decay or integrability conditions.
Their Section~5.3.1, pp.~443--444, sketches quadratic and modular
extensions. These interpolation identities do not identify the growing
orthogonal theta projection or its harmonic-measure limit.

\item[Other Hilbert-space trace constructions.]
Li's 2012 abstract describes a trace on a sum of Sonine, prolate,
and related spaces. The exact theorem and cutoff were not
available for comparison here \citep{li2012}. No claim of
exhaustive literature coverage or priority is made. The
comparison above records the statements that enter the proof.
\end{description}
